\section{Introduction and statements}
A scheme defined by failure of a multiplication map can carry geometric
information which disappears upon reduction. We exhibit such information
in a finite local algebra: the scheme remembers a polarized K3 surface,
although its reduced divisor is a fixed rational Schubert variety.
The reconstruction is obtained from a line bundle intrinsic to the
embedded nilpotent structure, rather than from an ambient marking or
from birational cancellation.

Let $V$ be a complex vector space of dimension $e\in\{3,4\}$, and put
$p=e(e-1)/2$. For an $e$-dimensional relation space
$R\subset\Sym^2V$, write
\begin{equation}\label{eq:algebra}
 S_R=\Sym^2V/R,\qquad
 B_R=\C\oplus V\oplus S_R.
\end{equation}
Multiplication on $V$ is the quotient map $\gamma_R$, and
$(V\oplus S_R)S_R=0$. Thus the maximal ideal has cube zero.
Set $X_R=\Gr(e,V\oplus S_R)$ and let $\mathcal K$ be its
rank-$e$ tautological subbundle. A point $K$ determines the unital
subspace $\C1\oplus K\subset B_R$. For any $m\ge2$ define
\begin{equation}\label{eq:full-failure}
 \widehat D_R=V\!\left(\Fitt_0\coker
 [\Sym^m(\OO\oplus\mathcal K)\longrightarrow B_R\otimes\OO]\right).
\end{equation}
The scheme is independent of $m\ge2$; in particular, ordinary powers
of its ideal below must not be confused with changing $m$.

We regard a relation $q\in R$ as a quadratic polynomial on $V^*$.
Let $G_e^\circ\subset\Gr(e,\Sym^2V)$ consist of the basepoint-free
systems whose Jacobian hypersurface
\begin{equation}\label{eq:Y-intro}
 Y_R=V\!\left(\det\left(\partial q_i/\partial x_j\right)\right)
       \subset\Pj(V^*)
\end{equation}
is smooth. This is a nonempty invariant open set. Its members are
ramification schemes of finite quadratic morphisms of degree
$2^{e-1}$; $Y_R$ is a genus-one curve for $e=3$ and a quartic K3
surface for $e=4$.

Write $\widehat\Delta_R=\{\rank(K\to V)<e\}$, and let $U_R$
be the open set where this rank is at least $e-1$. Set
$D_R=\widehat D_R\cap U_R$ and $\Delta_R=\widehat\Delta_R\cap U_R$.
If $\NN_R$ is the nilradical of $\OO_{D_R}$, its annihilator defines
an intrinsic smooth scheme $E_R\subset D_R$, and $\NN_R$ is an
invertible $\OO_{E_R}$-module. The following is the main reconstruction
statement.

\begin{theorem}[Polarized reconstruction]\label{thm:main-reconstruction}
For $R\in G_e^\circ$, put
\[
 \LL_R=\omega_{E_R}\otimes\NN_R^{\otimes-(p+e-1)}.
\]
There is an isomorphism of graded algebras
\begin{equation}\label{eq:main-ring}
 \bigoplus_{n\ge0}H^0(E_R,\LL_R^{\otimes n})
 \simeq
 \bigoplus_{n\ge0}H^0(Y_R,\OO_{Y_R}(en)).
\end{equation}
Consequently $D_R$ determines $(Y_R,\OO_{Y_R}(e))$.
For $e=4$ it determines $(Y_R,\OO_{Y_R}(1))$.
The same conclusions hold with $\widehat D_R$ in place of $D_R$.
No isomorphism is required to preserve $X_R$, $R$, or any projection.
\end{theorem}

The passage from $E_R$ to its section ring is essential. Stable
birational cancellation recovers the abstract surface, but not the
polarization. Here a cancellation of two determinant line bundles
produces the fourth power of the quartic polarization; its section
ring descends without extra functions in the affine-bundle directions.

\begin{theorem}[Variation and finite ambiguity]\label{thm:main-moduli}
For $e=4$, the image of
$R\mapsto(Y_R,\OO_{Y_R}(1))$ in the moduli space of quartic K3
surfaces has dimension nine. On a nonempty invariant open, the induced
map from projective equivalence classes of relation spaces to its image
is generically finite. Thus, over a general point of this image,
$D_R$ determines at most finitely many isomorphism classes of algebras
$B_R$; the same is true for $\widehat D_R$.
\end{theorem}

The dimension assertion is a statement about the actual ramification
surfaces, not an orbit count for unrelated algebras. We prove it by a
$25\times25$ integral determinant for the differential of the
Jacobian map on a $24$-dimensional Grassmannian chart.

The next statement describes the first higher-corank stratum. Its
precise open hypotheses are given in Definition~\ref{def:good-pair}.
They are verified by the same rational relation system used for the
moduli calculation.

\begin{theorem}[A corank-two embedded component]\label{thm:main-coranktwo}
There is a nonempty invariant open $G_4^\dagger\subset G_4^\circ$
such that, for each $R\in G_4^\dagger$, a dense open set of
rank-two images $H\subset V$ satisfies the following assertion.
Near any plane $K$ with image such an $H$, after an \emph{\'etale}
local extension and smooth auxiliary coordinates, the fixed-tensor
failure ideal has an irredundant primary decomposition
\begin{equation}\label{eq:main-coranktwo}
 \II_{\widehat D_R}
   =(\delta)\cap P_1^2\cap P_2^2\cap P_3^2\cap P_4^2\cap\mathfrak m^5,
 \qquad \delta=ad-bc,\quad\mathfrak m=(a,b,c,d).
\end{equation}
Each $P_i$ is the ideal of matrices whose image lies in one of the four
ramification lines in $V/H$. Besides the reduced divisor and the four
local ramification branches, the rank-two stratum is an associated
support. Its nilradical has index three. The square of the nilradical
is an invertible module on that stratum, of transverse length one.
\end{theorem}

In particular, the higher-corank contribution is not obtained by
continuing the two-parameter model $t(t,f)$ unchanged: it contains
additional associated information. The theorem is a fixed-tensor
calculation, not a specialization of a generic-matrix primary
intersection. It covers the stated open part of corank two; it does
not classify its exceptional complement or projection corank at least
three. By contrast, Theorem~\ref{thm:main-reconstruction} concerns
isomorphisms of the failure scheme on the whole $X_R$.

\subsection*{Relation with classical constructions}
The quadratic morphism and its bilinear incidence surface are classical
linear-system constructions. In particular, the incidence realization
of the K3 cover associated with a web of quadrics is recalled in
\cite[\S4]{IK}. Our contribution is the identification and use of its
polarization \emph{inside the intrinsic nilpotent structure of a
multiplication-failure scheme}, together with the fixed-tensor
corank-two calculation. We give the direct incidence comparison in
\S\ref{sec:classical}. The weighted determinantal formula retained in
Appendix~\ref{sec:weighted} is derived coefficient by coefficient from
the classical symbolic-order profile; it is not proposed as a replacement
for the De Concini--Eisenbud--Procesi theory.

The historical failure-locus comparison with Ballico's 1993 article
\cite{Ballico93} has not been completed at theorem-text level. Neither
metadata nor the later failure-cycle paper can settle that comparison.
This outstanding documentary issue is recorded separately from the
mathematical statements proved here. The complete companion preserves
the preceding manuscript, including its conductor, global-section,
classificatory, and application developments, without deleting them.
